%% Engineering thesis: main file, built on the WIT PWr thesis template
%% (thesis_templates/szablonwitpd, 25.02.2026). Its layout follows the template's main.tex.
%%
%% Build:  latexmk -pdf thesis.tex
%%   or:   pdflatex thesis && bibtex thesis && pdflatex thesis && pdflatex thesis
%%
%% Files: settings.tex (the faculty's page layout, fonts, title page and macros),
%% settings-en.tex (the English version of the template: title page, PDF metadata,
%% abstracts, \cref names, page numbers), abstract.tex, acronyms.tex, chapters/,
%% figures/, references.bib (exported by Zotero + Better BibTeX).

\documentclass[12pt, a4paper, twoside, titlepage]{report}

% Options for packages that settings.tex loads (they must come before it)
\PassOptionsToPackage{main=english}{babel}           % English text, Polish for the Streszczenie
\PassOptionsToPackage{capitalise,noabbrev}{cleveref} % \cref{fig:x} -> Figure 3.1

% Bibliography: loaded before settings.tex, so that hyperref (loaded there) links the citations
\RequirePackage[numbers,sort&compress]{natbib} % \citep{key} -> [1], \citet{key} -> Wang et al. [1]

%%% Copyright (C) [2026] [PWr * WIT * Wydziałowa Komisja d/s Jakości Kształcenia]
%%% Zmodyfikowany i zoptymalizowany (2026)
%%% ---------------------------------------------------------------------------
%%% [thesis] Faculty template file (thesis_templates/szablonwitpd, 25.02.2026).
%%% Changes are marked [thesis]: fallbacks for packages missing from this TeX
%%% install. The English version of the template's macros is in settings-en.tex.

%%% KONFIGURACJA EKSTRAKCJI TEKSTU (PDF COPY-PASTE)
\input{glyphtounicode} 
\pdfgentounicode=1     % Wymagane przez JSA (Jednolity System Antyplagiatowy)

%%% 1. KODOWANIE I JĘZYK
\RequirePackage[T1]{fontenc}
%\RequirePackage[utf8]{inputenc} % Opcjonalne w nowych systemach (domyślne)
\RequirePackage[english,polish]{babel} % Kolejność eng -> pl
\RequirePackage{csquotes} % Wymagane przez biblatex

%%% 2. MATEMATYKA I CZCIONKI
\RequirePackage{amsmath}
\RequirePackage{mathtools}
\RequirePackage{amsthm}

%%% [thesis] newtx is in Ubuntu's texlive-fonts-extra. Without it, fall back to the
%%% TeX Gyre Times and Helvetica clones (text) and mathptmx (maths), so the thesis still
%%% compiles; once newtx is installed it is used automatically, as in the template.
%%% (TeX Gyre, not psnfss Times: psnfss builds ą/ę from two glyphs, which breaks
%%% copying Polish text out of the PDF, and the JSA check reads that text.)
\IfFileExists{newtxtext.sty}{%
  \RequirePackage{newtxtext}
  \RequirePackage{newtxmath}
}{%
  \RequirePackage{mathptmx}
  \RequirePackage{tgtermes}
  \RequirePackage[scale=0.92]{tgheros}
}
\RequirePackage{microtype} % Poprawia optyczny wygląd tekstu

%%% 3. GEOMETRIA I GRAFIKA
\RequirePackage{geometry}
\geometry{
  a4paper,
  top=25mm, bottom=25mm,
  inner=25mm, outer=25mm,
  bindingoffset=5mm
}
\RequirePackage{graphicx}
\RequirePackage{xcolor}
\RequirePackage{float} % Dla opcji pozycjonowania [H]

%%% 4. NARZĘDZIA POMOCNICZE
\RequirePackage{etoolbox}

%%% 5. FORMATOWANIE TEKSTU I NAGŁÓWKÓW
\RequirePackage{titlesec}
\RequirePackage{caption}
\RequirePackage{subcaption}
\RequirePackage{setspace}

\setlength{\parindent}{0.7cm}
\setstretch{1.15} 

% Formatowanie nagłówków (Rozdział, Podrozdział)
\titleformat{\chapter}[display]
  {\bfseries\Large}
  {\chaptertitlename\ \thechapter}{0.5ex}   
  {\huge}
\titlespacing*{\chapter}{0pt}{-20pt}{20pt}

\titleformat{\section}{\bfseries\Large}{\thesection.}{1em}{}
\titlespacing*{\section}{0pt}{2.5ex plus 1ex minus .2ex}{1.3ex plus .2ex}

\titleformat{\subsection}{\bfseries\large}{\thesubsection.}{1em}{}
\titlespacing*{\subsection}{0pt}{2.25ex plus 1ex minus .2ex}{1.0ex plus .2ex}

% Numeracja ciągła wewnątrz rozdziałów
\counterwithin{figure}{chapter}
\counterwithin{table}{chapter}
\numberwithin{equation}{chapter}

% Podpisy (Captions) - pogrubione etykiety
\captionsetup{font=small, labelfont=bf, labelsep=period}

%%% 6. ALGORYTMY I LISTINGI KODU
%%% [thesis] algorithm/algpseudocode are in Ubuntu's texlive-science; load them only
%%% when installed (the thesis has no pseudocode yet).
\newif\ifstAlgorithms
\IfFileExists{algpseudocode.sty}{\stAlgorithmstrue}{}
\ifstAlgorithms
\RequirePackage{algorithm}
\RequirePackage{algpseudocode}
\algrenewcommand\algorithmiccomment[1]{\hfill\(\triangleright\) #1}
\floatname{algorithm}{Algorytm}
\fi

\RequirePackage{listings}
\definecolor{codebg}{RGB}{250,250,250}
\definecolor{codekw}{RGB}{0,64,160}
\definecolor{codestr}{RGB}{0,128,0}
\definecolor{codecmt}{RGB}{128,128,128}

\lstset{
  basicstyle=\ttfamily\small,
  backgroundcolor=\color{codebg},
  keywordstyle=\bfseries\color{codekw},
  stringstyle=\color{codestr},
  commentstyle=\itshape\color{codecmt},
  numbers=left,
  numberstyle=\tiny\color{gray},
  stepnumber=1,
  numbersep=8pt,
  tabsize=2,
  showstringspaces=false,
  breaklines=true,
  frame=lines,
  captionpos=b,
  % Mapowanie polskich znaków dla listings (niezbędne przy tym pakiecie)
  literate={ą}{{\k{a}}}1 {ć}{{\'c}}1 {ę}{{\k{e}}}1 {ł}{{\l{}}}1 {ń}{{\'n}}1 {ó}{{\'o}}1 {ś}{{\'s}}1 {ż}{{\.z}}1 {ź}{{\'z}}1
           {Ą}{{\k{A}}}1 {Ć}{{\'C}}1 {Ę}{{\k{E}}}1 {Ł}{{\L{}}}1 {Ń}{{\'N}}1 {Ó}{{\'O}}1 {Ś}{{\'S}}1 {Ż}{{\.Z}}1 {Ź}{{\'Z}}1
}
\renewcommand{\lstlistingname}{Listing}
\renewcommand{\lstlistlistingname}{Spis listingów}

% Styl Haskell 
\definecolor{darkgreen}{rgb}{0,0.5,0}
\lstdefinestyle{haskellStyle}{
    language=Haskell,
    commentstyle=\color{darkgreen},
    keywordstyle=\color{blue}\bfseries,
    basicstyle=\ttfamily\small,
    breaklines=true,
    stringstyle=\color{red}
}

%%% 7. HYPERREF I CLEVEREF (Musi być na końcu)
\RequirePackage[
    pdfusetitle,
    colorlinks=true,
    linkcolor=black,
    citecolor=blue,
    urlcolor=blue,
    bookmarksnumbered=true
]{hyperref}

\RequirePackage[nameinlink]{cleveref}

% Tłumaczenia dla cleveref
\crefname{figure}{Rys.}{Rys.}
\Crefname{figure}{Rysunek}{Rysunki}
\crefname{table}{Tab.}{Tab.}
\Crefname{table}{Tabela}{Tabele}
\crefname{equation}{równ.}{równ.}
\Crefname{equation}{Równanie}{Równania}
\crefname{listing}{listing}{Listingi}
\crefname{algorithm}{algorytm}{algorytmy}
\Crefname{algorithm}{Algorytm}{Algorytmy}

%%%%%%%%%%%
%% MAKRA %%
%%%%%%%%%%%

% Definicja rodzaju pracy 
\newcommand{\stRodzajPracy}{%
  \ifnum\stCzyMgr=1 MAGISTERSKA\else INŻYNIERSKA\fi
}

\newcommand{\StronaTytulowa}{
  \pagenumbering{roman}
  \sffamily
  \begin{titlepage}  
    \newgeometry{top=2.5cm, bottom=2.5cm, left=2.5cm, right=2.5cm}
    
    \begin{center}
      {\huge\bfseries Politechnika Wrocławska \par}
      \vspace{0.4cm}
      {\Large Wydział Informatyki i Telekomunikacji \par}
    \end{center}
    \par\noindent\rule{\textwidth}{1pt}\par 
    \vspace{0.5cm} 
    
    {\large
    \begin{tabular}{@{}ll}
      Kierunek: & \textbf{\stKierunek\ (\stKodKierunku)} \\
      \ifdefempty{\stSpecjalnosc}{}{Specjalność: & \textbf{\stSpecjalnosc\ (\stKodSpecjalnosci)}}
    \end{tabular}
    }
    
    \vspace*{\fill}
    
    \begin{center}
      {\huge\bfseries PRACA DYPLOMOWA \par}
      \vspace{0.3cm}    
      {\huge \stRodzajPracy \par}
      \vspace{2.0cm}
      {\Large \bfseries \stTytulPL \par} 
      \vspace{1cm}
      {\Large \stAutor \par}
      \vspace{2cm}
      {\large Opiekun pracy:\par \textbf{\stOpiekun} \par}
    \end{center}
    
    \vspace*{\fill}
    
    {\raggedright
      \small Słowa kluczowe: \textit{\stSlowaKluczowe}\par
      \vspace{0.5em}
      \par\noindent\rule{\textwidth}{1pt}\par 
    }
    
    \begin{center}
      \Large Wrocław \stRok
    \end{center}
  
    \restoregeometry
  \end{titlepage}
  \normalfont
  \pagestyle{empty}
}

\newenvironment{abstracts}{}{\selectlanguage{polish}\clearpage}

\newcommand{\AbstractPL}{
  \begin{center}{\Large\bfseries Streszczenie \par}\end{center}\vspace{0.5em}
}

\newcommand{\AbstractENG}{
  \selectlanguage{english}
  \begin{center}{\Large\bfseries Abstract \par}\end{center}\vspace{0.5em}
}

\newcommand{\SpisTresci}{
  \clearpage
  \addtocontents{toc}{\protect\thispagestyle{empty}} 
  \pagestyle{empty} 
  \tableofcontents
  \clearpage
  \pagestyle{empty}
}

%% WSTAWIENIE METADANYCH DO PLIKU PDF
\newcommand{\SetMetaData}{%
  \AtBeginDocument{%  
    \hypersetup{%
      pdfinfo = {%
        Title    = {\stTytulPL},
        Author   = {\stAutor},
        Subject  = {Praca dyplomowa \stRodzajPracy},  
        Keywords = {\stSlowaKluczowe}, 
        ModDate  = {\pdfcreationdate}, 
        Creator  = {pdftex}
      }
    }
  }
}
%%% KONIEC PLIKU SETTINGS.TEX    % the faculty template: page layout, fonts and macros
%%% English version of the faculty template, loaded by thesis.tex right after settings.tex.
%%% settings.tex is the faculty's file and stays as it is. Everything that differs because
%%% this thesis is written in English is redefined here, plus page numbers on every page.
%%% The thesis data (title, author, keywords, ...) is in thesis.tex.

%% 1. TITLE PAGE
% The template's \StronaTytulowa (settings.tex) with the same layout and English text.
% Changes: the labels and names are in English, the title is \stTytulAng, and the page style
% after the title page is plain (page numbers) instead of empty.
\newcommand{\stRodzajPracyAng}{%
  \ifnum\stCzyMgr=1 MASTER'S THESIS\else ENGINEERING THESIS\fi
}

\renewcommand{\StronaTytulowa}{
  \pagenumbering{roman}
  \sffamily
  \begin{titlepage}
    \newgeometry{top=2.5cm, bottom=2.5cm, left=2.5cm, right=2.5cm}

    \begin{center}
      {\LARGE\bfseries Wrocław University of Science and Technology \par}
      \vspace{0.4cm}
      {\Large Faculty of Information and Communication Technology \par}
    \end{center}
    \par\noindent\rule{\textwidth}{1pt}\par
    \vspace{0.5cm}

    {\large
    \begin{tabular}{@{}ll}
      Field of study: & \textbf{\stKierunek\ (\stKodKierunku)} \\
      \ifdefempty{\stSpecjalnosc}{}{Specialization: & \textbf{\stSpecjalnosc\ (\stKodSpecjalnosci)}}
    \end{tabular}
    }

    \vspace*{\fill}

    \begin{center}
      {\huge\bfseries \stRodzajPracyAng \par}
      \vspace{2.0cm}
      {\Large \bfseries \stTytulAng \par}
      \vspace{1cm}
      {\Large \stAutor \par}
      \vspace{2cm}
      {\large Supervisor:\par \textbf{\stOpiekun} \par}
    \end{center}

    \vspace*{\fill}

    {\raggedright
      \small Keywords: \textit{\stSlowaKluczowe}\par
      \vspace{0.5em}
      \par\noindent\rule{\textwidth}{1pt}\par
    }

    \begin{center}
      \Large Wrocław \stRok
    \end{center}

    \restoregeometry
  \end{titlepage}
  \normalfont
  \pagestyle{plain} % the template has \pagestyle{empty}: no page numbers
}

%% 2. PDF METADATA (file properties): English title, thesis type and keywords
\renewcommand{\SetMetaData}{%
  \AtBeginDocument{%
    \hypersetup{%
      pdfinfo = {%
        Title    = {\stTytulAng},
        Author   = {\stAutor},
        Subject  = {\ifnum\stCzyMgr=1 Master's thesis\else Engineering thesis\fi},
        Keywords = {\stSlowaKluczowe},
        ModDate  = {\pdfcreationdate},
        Creator  = {pdftex}
      }
    }
  }
}

%% 3. ABSTRACTS
% The template assumes Polish is the main language. Here English is (the babel option in
% thesis.tex), so each abstract selects its own language and the text returns to English
% after them. \AbstractENG from settings.tex already selects English.
\renewenvironment{abstracts}{}{\selectlanguage{english}\clearpage}
\renewcommand{\AbstractPL}{%
  \selectlanguage{polish}%
  \begin{center}{\Large\bfseries Streszczenie \par}\end{center}\vspace{0.5em}%
}

%% 4. PAGE NUMBERS ON EVERY PAGE EXCEPT THE TITLE PAGE
% The template's \StronaTytulowa and \SpisTresci switch page numbers off, and nothing switches
% them back on, so only chapter-opening pages had numbers. Here every page after the title
% page has one: roman (ii, iii, ...) for the abstracts, contents and abbreviations, arabic
% from Chapter 1 on (thesis.tex switches the numbering there, as in the template).
\renewcommand{\SpisTresci}{\clearpage\tableofcontents\clearpage}

%% 5. NAMES FOR \cref (settings.tex sets Polish ones)
\crefname{figure}{Figure}{Figures}
\Crefname{figure}{Figure}{Figures}
\crefname{table}{Table}{Tables}
\Crefname{table}{Table}{Tables}
\crefname{equation}{Equation}{Equations}
\Crefname{equation}{Equation}{Equations}
\crefname{listing}{Listing}{Listings}
\Crefname{listing}{Listing}{Listings}
\crefname{algorithm}{Algorithm}{Algorithms}
\Crefname{algorithm}{Algorithm}{Algorithms}
\ifstAlgorithms \floatname{algorithm}{Algorithm} \fi
 % its English version: title page, PDF metadata, abstracts, \cref names

%% THESIS DATA (printed on the title page and in the PDF metadata)

\newcommand{\stKierunek}{Algorithmic Computer Science}
\newcommand{\stKodKierunku}{INA}
\newcommand{\stSpecjalnosc}{} % leave empty if there is no specialization
\newcommand{\stKodSpecjalnosci}{} % the codes of fields and specializations are on the WIT pages
\newcommand{\stCzyMgr}{0} % 1 = master's thesis, 0 = engineering thesis
\newcommand{\stTytulAng}{Design and Implementation of a Real-Time Brain–Computer Interface for
  Predictive EEG-Based Communication Using Machine Learning and Natural Language Processing}
\newcommand{\stTytulPL}{Projekt i implementacja interfejsu mózg–komputer działającego w czasie
  rzeczywistym do predykcyjnej komunikacji opartej na EEG z wykorzystaniem uczenia maszynowego
  i przetwarzania języka naturalnego} % not printed: the title page is in English
\newcommand{\stAutor}{Michał Szandarowski}
\newcommand{\stOpiekun}{dr inż. Michał Kędziora}
\newcommand{\stSlowaKluczowe}{brain–computer interface, BCI, electroencephalography, EEG,
  steady-state visual evoked potentials, SSVEP, real-time signal processing, spatial
  filtering, CCA, TRCA, machine learning, natural language processing, text prediction,
  information transfer rate, ITR, assistive communication}
\newcommand{\stRok}{2026} % year of submission

\SetMetaData % PDF metadata (title, author, keywords)

%% PACKAGES USED BY THE THESIS (on top of settings.tex)
\IfPackageLoadedTF{newtxmath}{}{\usepackage{amssymb}} % newtxmath already defines these symbols
\usepackage{bm}
\graphicspath{{figures/}}
\usepackage{booktabs,tabularx,multirow}
\usepackage{tikz}
\usetikzlibrary{arrows.meta,positioning,shapes.geometric,fit}
\usepackage{enumitem}
\usepackage{xurl}

%% OBSŁUGA BIBLIOGRAFII (bibliography: natbib is loaded at the top of this file)
\bibliographystyle{unsrtnat} % numbered in order of first citation, like the template's sorting=none

%% AKRONIMY (skróty;synonimy)
%%% List of abbreviations, with the acro package as in the faculty template.
%%% Use \ac{KEY} in the text: the first use prints "long form (SHORT)", later uses print
%%% "SHORT". \acp{KEY} = plural. \acs*{KEY} = short form only, without counting as the first
%%% use (in acro, plain \acs{KEY} counts as the first use, unlike in the old acronym package).
%%% Only the abbreviations used in the text are printed, sorted by their short form.
%%% Compile a few times so the list and its page numbers settle.

\RequirePackage{acro}

% Width of the short-form column in the list (acro's lof style takes the list-of-figures
% layout, 2.3em wide, which is too narrow for FB-eTRCA or SSVEP).
\makeatletter
\renewcommand*{\l@acro}{\@dottedtocline{1}{1.5em}{5.5em}}
\makeatother

%% Uważaj na formatowanie;
%% Nie zapomnij o przecinkach

\DeclareAcronym{ALS}{short = ALS, long = amyotrophic lateral sclerosis}
\DeclareAcronym{BCI}{short = BCI, long = brain--computer interface}
\DeclareAcronym{CCA}{short = CCA, long = canonical correlation analysis}
\DeclareAcronym{CI}{short = CI, long = continuous integration}
\DeclareAcronym{EEG}{short = EEG, long = electroencephalography}
\DeclareAcronym{eTRCA}{short = eTRCA, long = ensemble task-related component analysis}
\DeclareAcronym{FBCCA}{short = FBCCA, long = filter bank canonical correlation analysis}
\DeclareAcronym{FBeTRCA}{short = FB-eTRCA, long = filter bank ensemble task-related component analysis}
\DeclareAcronym{GUI}{short = GUI, long = graphical user interface}
\DeclareAcronym{GPU}{short = GPU, long = graphics processing unit}
\DeclareAcronym{ITR}{short = ITR, long = information transfer rate}
\DeclareAcronym{JSON}{short = JSON, long = JavaScript Object Notation}
\DeclareAcronym{LDA}{short = LDA, long = linear discriminant analysis}
\DeclareAcronym{LM}{short = LM, long = language model}
\DeclareAcronym{LSL}{short = LSL, long = Lab Streaming Layer}
\DeclareAcronym{NIRS}{short = NIRS, long = near-infrared spectroscopy}
\DeclareAcronym{NLP}{short = NLP, long = natural language processing}
\DeclareAcronym{PSD}{short = PSD, long = power spectral density}
\DeclareAcronym{SNR}{short = SNR, long = signal-to-noise ratio}
\DeclareAcronym{SSVEP}{short = SSVEP, long = steady-state visual evoked potential}
\DeclareAcronym{TRCA}{short = TRCA, long = task-related component analysis}

%%%%%%%%%%%%%%%%%%%%%%%%%%%%%%%%%%%%%%%%%%%
%% NIE RUSZAJ TEGO CO JEST NIŻEJ JEŚLI  %%%
%% NIE ZNASZ ZBYT DOBRZE LATEX'A        %%%
%%%%%%%%%%%%%%%%%%%%%%%%%%%%%%%%%%%%%%%%%%%
%% [thesis] Only change from the template: the heading is in English.
\newcommand{\ShowAcronyms}[1]{
  \begin{center}
    {\Large\bfseries List of abbreviations \par}
  \end{center}
  \vspace{0.5em}

  \ifboolexpr{
      test {\ifstrequal{#1}{lof}} or
      test {\ifstrequal{#1}{toc}} or
      test {\ifstrequal{#1}{description}}
    }{%
      \printacronyms[sort=true, heading=none, template=#1]
    }{%
      \printacronyms[sort=true, heading=none, template=lof]
     }%
  \clearpage
}


%% DRAFT HELPERS (switch todonotes to [disable] for the final version)
\usepackage[colorinlistoftodos]{todonotes}
\newcommand{\missing}[1]{\todo[inline,color=red!15]{Missing: #1}}

% Build only the chapters you are working on (faster). Comment out for the full thesis.
% \includeonly{chapters/03_system}

%% TU ZACZYNA SIĘ DOKUMENT

\begin{document}

%% NIE RUSZAJ NASTĘPNYCH TRZECH LINIJEK %%
\StronaTytulowa
% Abstracts in English and Polish. Write them last.
% Outline: a one-page summary: the problem, the approach (an SSVEP speller on a low-cost
% dry-electrode cap, supported by a language model) and the main results.
% Don't use \ac here: the abstracts define their own abbreviations in plain text.
\begin{abstracts}

  \AbstractENG

  This thesis concerns the design and implementation of a non-invasive, real-time
  brain–computer interface (BCI) intended to support communication for people with severe
  motor impairments. The system is based on the steady-state visual evoked potential
  (SSVEP) paradigm. The user selects characters or words by focusing their gaze on
  interface elements that flicker at different frequencies.

  The engineering part includes several components: a graphical user interface with
  precise visual stimulus timing and a signal processing pipeline.

  The research part compares spatial filtering and classification methods, including
  canonical correlation analysis (CCA) and task-related component analysis (TRCA),
  together with selected machine learning algorithms. The algorithms were trained and
  evaluated on offline recorded data. The system is integrated with a natural language
  processing (NLP) model that predicts subsequent words in real time. This is intended to
  increase communication speed and reduce user fatigue.

  The integrated system was validated in online experiments. Classification accuracy and
  information transfer rate (ITR) were measured both with and without text prediction. The
  results show to what extent combining SSVEP classification with language prediction
  improves communication efficiency in BCI systems.

  \AbstractPL

  Praca dotyczy projektu i implementacji nieinwazyjnego interfejsu mózg–komputer (BCI)
  działającego w czasie rzeczywistym, przeznaczonego do komunikacji osób z ciężkimi
  dysfunkcjami ruchowymi. System wykorzystuje paradygmat potencjałów wywołanych stanu
  ustalonego (SSVEP). Użytkownik wybiera w nim znaki lub słowa, skupiając wzrok na
  elementach interfejsu migających z różnymi częstotliwościami.

  Część inżynierska obejmuje kilka elementów: graficzny interfejs użytkownika z precyzyjną
  synchronizacją bodźców wzrokowych oraz potok przetwarzania sygnału.

  Część badawcza polega na porównaniu metod filtracji przestrzennej i klasyfikacji, m.in.
  analizy korelacji kanonicznej (CCA) i analizy komponentów powiązanych z zadaniem (TRCA),
  a także wybranych algorytmów uczenia maszynowego. Algorytmy wytrenowano i oceniono na
  danych zarejestrowanych offline. System zintegrowano z modelem przetwarzania języka
  naturalnego (NLP), który w czasie rzeczywistym podpowiada kolejne słowa. Ma to zwiększyć
  szybkość komunikacji i ograniczyć zmęczenie użytkownika.

  Działanie zintegrowanego systemu zweryfikowano w eksperymentach online. Oceniono
  dokładność klasyfikacji oraz szybkość transferu informacji (ITR) zarówno z predykcją
  tekstu, jak i bez niej. Wyniki pozwalają określić, w jakim stopniu połączenie
  klasyfikacji sygnałów SSVEP z predykcją językową poprawia efektywność komunikacji w
  systemach BCI.

\end{abstracts} % plik abstrakt.txt edytujesz samodzielnie
\SpisTresci

%%% Tutaj wyświetlamy listę akronimów
%%  jako parametr (template) możesz wybrać: lof, toc, description
%%%
\ShowAcronyms{lof}

\pagenumbering{arabic}  %% NIE RUSZAJ TEGO

%% TU JEST MIEJSCE NA DOŁĄCZENIE TWOICH ROZDZIAŁÓW
% Chapter 1. Write it near the end, after Chapters 3--7.
% Outline: why the work matters: people with ALS or locked-in syndrome keep their minds
% but lose muscle control, and a brain--computer interface can let them spell. The best
% SSVEP spellers are laboratory systems (64-channel gel EEG, 40 targets, no use of
% language). The chapter states the aim and scope of the thesis, its contributions and
% the structure of the text.

\chapter{Introduction}
\label{ch:introduction}

\section{Motivation}
\label{sec:motivation}

\section{Problem statement}
\label{sec:problem}

\section{Aim and scope}
\label{sec:aim}

\section{Contributions}
\label{sec:contributions}

\section{Structure of the thesis}
\label{sec:structure}

% Chapter 2 (thesis task 1).
% Outline: the theory and prior work the thesis builds on: what EEG measures and what
% disturbs it, the main BCI paradigms and why SSVEP was chosen, how SSVEP stimuli are
% generated and detected (CCA, FBCCA, TRCA), the best existing SSVEP spellers, language
% models in BCI spellers, and how performance is measured (accuracy, ITR). Ends with the
% research gap.

\chapter{Background and literature review}
\label{ch:background}

\section{EEG basics}
\label{sec:eeg}

\section{BCI paradigms}
\label{sec:paradigms}

\section{SSVEP and stimulus design}
\label{sec:ssvep}

\section{Detection methods}
\label{sec:detection}

\section{SSVEP spellers}
\label{sec:spellers}

\section{Language models in BCI}
\label{sec:lm-bci}

\section{Performance measures}
\label{sec:measures}

\section{Research gap}
\label{sec:gap}

% Chapter 3 (thesis tasks 2 and 3, part of 7).
% Outline: the engineering part: requirements, hardware (BrainAccess cap with dry
% electrodes, 120 Hz monitor), the two-process architecture connected by Lab Streaming
% Layer, the flicker stimulus and its frame timing, the 6-target two-level speller
% layout, the calibration protocol, EEG acquisition and recording, the online speller
% loop, and how the system was tested without the cap.

\chapter{System design and implementation}
\label{ch:system}

\section{Requirements}
\label{sec:requirements}

The system is intended for people with \ac{ALS} or locked-in syndrome, who have
lost most muscle control but retain their cognitive function and can still move their
eyes (\cref{sec:motivation}).
It must work on low-cost hardware: an \ac{EEG} cap with a low number of channels and
dry electrodes, which need no gel.
With it, the user must be able to spell with high accuracy and speed, helped by a
language model that predicts the next character or word.
\Cref{tab:requirements} lists the functional requirements, which describe what the
system does, and the non-functional requirements, which constrain how it does it.

\begin{table}[tb]
  \centering
  \caption{Functional (FR) and non-functional (NFR) requirements.}
  \label{tab:requirements}
  \begin{tabularx}{\textwidth}{@{}lXl@{}}
    \toprule
    ID & Requirement: the system must\ldots & Verified in \\
    \midrule
    \multicolumn{3}{@{}l}{\textit{Functional requirements}} \\
    FR1 & present six targets flickering at chosen frequencies between 8 and 15\,Hz, not limited to divisors of the refresh rate & \\
    FR2 & timestamp every stimulus event on the same clock as the \acs{EEG} & \\
    FR3 & run a calibration session that records labelled \acs{EEG} trials for training the classifier & \\
    FR4 & record the \acs{EEG} and the stimulus events to a file for offline analysis & \\
    FR5 & classify each selection in real time and show the selected target to the user & \\
    FR6 & let the user type any English text with only six targets: the 26 letters, space and backspace, a way to undo a wrong selection, and accepting a suggested word & \\
    FR7 & suggest words from a language model that can be switched on and off, so that the online experiment can compare both conditions & \\
    FR8 & support copy-spelling of a given phrase: cue the next correct target and log it as the ground truth for accuracy and \acs{ITR} & \\
    FR9 & replay a recording offline through the online processing pipeline & \\
    \midrule
    \multicolumn{3}{@{}l}{\textit{Non-functional requirements}} \\
    NFR1 & keep the stimulus timing precise, with no dropped frames in normal operation & \\
    NFR2 & finish the classification and the word suggestions within the 0.5\,s gaze shift & \\
    NFR3 & process and classify the \acs{EEG} identically offline and online & \\
    NFR4 & be testable end to end without the \acs{EEG} cap & \\
    NFR5 & run on a consumer laptop without a dedicated graphics card, with a low-cost dry-electrode \acs{EEG} cap & \\
    NFR6 & be publicly available, with automated tests that run on every commit (continuous integration) & \\
    NFR7 & be used only after the user has been warned about the risk of photosensitive epilepsy & \\
    \bottomrule
  \end{tabularx}
\end{table}

\section{Hardware}
\label{sec:hardware}

\section{Architecture}
\label{sec:architecture}

% (moved from 3.1) The system should be modular, with separate processes for stimulus
% presentation and EEG acquisition, and communicate via Lab Streaming Layer.

\section{Stimulus presentation and timing}
\label{sec:stimulus}

\section{Speller layout}
\label{sec:layout}

\section{Calibration protocol}
\label{sec:calibration}

\section{EEG acquisition and recording}
\label{sec:acquisition}

\section{Online speller}
\label{sec:online}

\section{Testing and verification}
\label{sec:testing}

% Chapter 4 (thesis tasks 4 and 5).
% Outline: the data and the methods: the Benchmark dataset and own recordings, the
% preprocessing pipeline (notch, band-pass, filter bank), the classifiers (CCA, FBCCA,
% TRCA, eTRCA, FB-eTRCA, and a PSD+LDA baseline), their implementation, and the
% evaluation method (cross-validation, accuracy and ITR against window length,
% statistical tests).

\chapter{Signal processing and classification}
\label{ch:processing}

\section{Datasets}
\label{sec:datasets}

\section{Preprocessing}
\label{sec:preprocessing}

\section{CCA and FBCCA}
\label{sec:cca}

\section{The TRCA family}
\label{sec:trca}

\section{Baseline}
\label{sec:baseline}

\section{Evaluation method}
\label{sec:evaluation}

% Chapter 5 (thesis task 6).
% Outline: how word prediction makes up for the small number of targets: word
% completion with a word-frequency model, context-aware re-ranking with a small
% transformer (distilgpt2), and combining the classifier's output with the language
% model's prediction. Includes the offline evaluation (keystroke savings, prediction
% latency).

\chapter{Language model integration}
\label{ch:language}

\section{Why a language model}
\label{sec:why-lm}

\section{Word completion}
\label{sec:completion}

\section{Transformer re-ranking}
\label{sec:reranking}

\section{Fusion with the classifier}
\label{sec:fusion}

\section{Offline evaluation}
\label{sec:lm-evaluation}

% Chapter 6 (thesis task 7).
% Outline: all results: the comparison of methods on the Benchmark dataset, stimulus
% timing measurements, results on own recordings and the choice of method and window,
% the online copy-spelling experiment protocol, and the online results with and without
% word prediction (accuracy, ITR, characters per minute).

\chapter{Experiments and results}
\label{ch:results}

\section{Benchmark comparison}
\label{sec:benchmark}

\section{Stimulus timing}
\label{sec:timing-results}

\section{Own recordings}
\label{sec:own-recordings}

\section{Online experiment}
\label{sec:online-experiment}

\section{Online results}
\label{sec:online-results}
% Numbers moved here from Ch. 3 (2026-11-23); Ch. 3 now points to this section for them.
% Nov 16 online session:
% - latency 0.6 s, chosen from a sweep over the calibration of the same session
% - answer delay (end of flicker -> answer): median 1.19 s, above the bound
%   tau + 0.5 s = 1.1 s, so the EEG timestamps did not lead the markers that day
% - classification: median 11 ms (max 22 ms); transformer suggestions whenever the
%   text changed: median 26 ms (max 74 ms), a small part of the wait
% - time per selection, onset to onset, used for the online ITR: 4.2 s =
%   0.5 s gaze shift + 2 s flicker + 1.2 s wait + 0.5 s feedback

\section{Language model results}
\label{sec:lm-results}

% Chapter 7.
% Outline: what the results mean and how they compare with the literature, the
% limitations (small group of healthy volunteers, gaze control required, English only),
% future work, and the final conclusions.

\chapter{Discussion and conclusions}
\label{ch:conclusions}

\section{Summary}
\label{sec:summary}

\section{Discussion}
\label{sec:discussion}

\section{Limitations}
\label{sec:limitations}

\section{Future work}
\label{sec:future}

\section{Conclusions}
\label{sec:conclusions}


%% TU WYWOŁUJEMY BIBLIOGRAFIĘ
\bibliography{references}

%% TU ZACZYNAJĄ SIĘ ZAŁĄCZNIKI
\appendix
% Appendices: each \chapter here becomes Appendix A, B, ...
% Outline: how to install and run the software, data formats and the communication
% protocol, the consent and epilepsy screening forms, and the phrases used in the
% experiment.

\chapter{Software installation and use}
\label{app:software}

\chapter{Data formats and communication protocol}
\label{app:protocol}

\chapter{Phrases used in the experiment}
\label{app:phrases}


%% --- Spisy ---
\listoffigures
\listoftables

\end{document}
